\documentclass[aps,prx,twocolumn,superscriptaddress,longbibliography,nofootinbib,floatfix]{revtex4-2}

\usepackage{amsmath,amssymb,amsthm,mathtools,bm}
\usepackage{booktabs,array}
\usepackage[colorlinks=true,linkcolor=blue,citecolor=blue,urlcolor=blue]{hyperref}
\usepackage{microtype}

\newcommand{\Id}{\mathbf{1}}
\newcommand{\Tr}{\operatorname{Tr}}
\newcommand{\tr}{\operatorname{tr}}
\newcommand{\dd}{\mathrm{d}}
\newcommand{\ii}{\mathrm{i}}
\newcommand{\Occ}{\mathrm{o}}
\newcommand{\Emp}{\mathrm{e}}

\hypersetup{
  pdftitle={Gaussian Trajectory Geometry: Reference Anisotropy, Occupied Frames, and Conditional Spectra},
  pdfauthor={Class-A fermionic adaptive-circuit project}
}

\begin{document}

\title{Gaussian Trajectory Geometry: Reference Anisotropy, Occupied Frames, and Conditional Spectra}
\date{August 29, 2026}

\begin{abstract}
We collect the Gaussian formalism behind three complementary numerical probes of an
adaptive class-A fermion circuit.  Pure trajectories are evolved as orthonormal occupied
frames, while maximally mixed trajectories are retained as correlation matrices.  For a
fixed nonzero-probability measurement record, differentiation of the normalized
correlation-matrix M\"obius map gives an exact congruence cocycle.  We prove that this
cocycle is the Jacobian of the occupied-frame evolution after quotienting the frame's
unitary gauge, and that its physical pure-state singular values are pairwise products of
occupied and empty one-leg singular values.  This identifies the invariant spectrum
targeted by the pure-state G5 estimator and distinguishes it from the trajectory-operator
spectrum reconstructed by the maximally mixed G4 calculation.  We also give the
charge-conserving Gaussian reference construction used to calibrate spacetime anisotropy
and work through the completed coarse-audit matching estimator without promoting it to a
final result.
\end{abstract}

\maketitle

\section{Gaussian coordinates and native occupied frames}

For complex fermions satisfying
$\{\hat\psi_i^\dagger,\hat\psi_j\}=\delta_{ij}$, the canonical correlation
matrix is
\begin{equation}
 G_{ij}=\Tr(\hat\rho\,\hat\psi_i^\dagger\hat\psi_j),
 \qquad 0\leq G\leq\Id .
 \label{eq:canonical-G}
\end{equation}
The implementation uses the transposed one-body convention
$C=G^{\mathsf T}$ and the centered covariance
\begin{equation}
 \Gamma=2C-\Id .
 \label{eq:centered}
\end{equation}
Transpose does not alter the spectra or Frobenius singular values discussed below.
A pure charge-conserving Gaussian state obeys $C^2=C$ and is represented by an
orthonormal occupied frame $F\in\mathbb C^{N\times r}$,
\begin{equation}
 \begin{gathered}
 F^\dagger F=\Id_r,\qquad C=FF^\dagger,\\
 F\sim FU,\qquad U\in U(r).
 \end{gathered}
 \label{eq:frame}
\end{equation}
Thus $F$ is a Stiefel representative of a point on $\operatorname{Gr}(r,N)$;
the right-unitary freedom is numerical gauge, not a physical degree of freedom.

Let $\chi$ be a normalized measured orbital, $P_\chi=\chi\chi^\dagger$, and
$Q_\chi=\Id-P_\chi$.  The occupied and empty Born probabilities are
$n_\chi=\chi^\dagger C\chi$ and $1-n_\chi$.  The two primitive rank-changing
frame operations have correlation matrices
\begin{align}
 C_{\rm gain}
 &=C+\frac{[(\Id-C)\chi][(\Id-C)\chi]^\dagger}{1-n_\chi},
 \label{eq:gain}\\
 C_{\rm loss}
 &=C-\frac{(C\chi)(C\chi)^\dagger}{n_\chi}.
 \label{eq:loss}
\end{align}
Numerically, gain appends the normalized residual $(\Id-C)\chi$ to $F$.
Loss computes $a=F^\dagger\chi$, applies a complex Householder reflector that
rotates $a$ to the last frame coordinate, and deletes that column.  An occupied
projector is loss followed by deterministic gain; an empty projector is gain
followed by deterministic loss.  These compositions give
\begin{align}
 C_{1}&=C-\frac{CP_\chi C}{n_\chi}+P_\chi,\nonumber\\
 C_{0}&=C+\frac{(\Id-C)P_\chi(\Id-C)}{1-n_\chi}-P_\chi.
 \label{eq:pure-projectors}
\end{align}
Thin QR is invoked only when the Gram residual warrants it.  The CPU engine stores one
variable-rank frame, whereas the GPU engine stores trajectory-dependent ranks in a padded
batch and applies the same append/Householder algebra selectively.  Frame-native trajectory
evolution therefore avoids repeatedly materializing an $N\times N$ covariance.  These
identities are the frame-level realization of charge-conserving Gaussian measurements
and feedforward~\cite{Bravyi2005,BhuiyanPanJian2026}.

\section{Normalized M\"obius tangent and the frame Jacobian}

Write the selected measurement outcome as $m\in\{0,1\}$,
$\eta_m=2m-1$, and the controller target as $s\in\{0,1\}$.  On the code-side
centered covariance, one normalized measurement followed by reset is
\begin{align}
 \Phi_{m|s}(\Gamma)
 &=Q_\chi(\Id+\eta_m\Gamma P_\chi)^{-1}\Gamma Q_\chi
   +\eta_sP_\chi,\nonumber\\
 p_m&=\frac{1+\eta_m g}{2},\qquad g=\tr(P_\chi\Gamma),
 \label{eq:mobius}
\end{align}
For $p_m>0$, the rank-one inverse is exact and
\begin{equation}
 L_m(\Gamma)=Q_\chi(\Id+\eta_m\Gamma P_\chi)^{-1}
 =Q_\chi-\frac{\eta_m Q_\chi\Gamma P_\chi}{1+\eta_m g}.
 \label{eq:L}
\end{equation}
Put $M=\Id+\eta_m\Gamma P_\chi$.  At fixed $P_\chi,m,s$, the identity
$\delta M^{-1}=-M^{-1}(\eta_mXP_\chi)M^{-1}$ gives
\begin{align}
 \delta\Phi
 &=Q_\chi M^{-1}X
 [\Id-\eta_mP_\chi M^{-1}\Gamma]Q_\chi\nonumber\\
 &=L_mXL_m^\dagger ,
 \label{eq:derivative-step}
\end{align}
because $[\Id-\eta_mP_\chi M^{-1}\Gamma]Q_\chi
=(\Id+\eta_mP_\chi\Gamma)^{-1}Q_\chi=L_m^\dagger$.  Thus
\begin{equation}
 \boxed{D_\Gamma\Phi_{m|s}[X]=L_mXL_m^\dagger.}
 \label{eq:event-tangent}
\end{equation}
The reset term is constant, so its target $s$ drops out.  Exact projectors really are
singular: $P_\chi L_m=0$ and, for pure $\Gamma$,
$\ker L_m=\operatorname{span}\{(\Id+\eta_m\Gamma)\chi\}$.  This vector lies in
the selected $\eta_m$ eigenspace: that one-leg block loses one rank, whereas the opposite
block is injective.  Such nulls are not pseudoinverted.  A unitary event
$\Gamma\mapsto U\Gamma U^\dagger$ has one-leg factor $L=U$.  Writing $L_t$ for either
kind of event factor, a fixed ordered record gives
\begin{align}
 A_{T:0}&=L_{T-1}\cdots L_1L_0,\nonumber\\
 X_T&=A_{T:0}X_0A_{T:0}^\dagger .
 \label{eq:cocycle}
\end{align}
This is a branchwise derivative of the normalized conditional state, not a derivative
through the outcome sampler and not a Born-score response.

We now make precise in what sense Eq.~\eqref{eq:cocycle} is the frame Jacobian.  Let
$f_{m|s}$ denote the exact gain/loss frame update and
$\pi_r(F)=2FF^\dagger-\Id$.  If $r'=r+s-m$, the state-level square commutes:
\begin{equation}
 \begin{array}{ccc}
 F&\xrightarrow{\ f_{m|s}\ }&F'\\[-1mm]
 {\scriptstyle\pi_r}\downarrow&&\downarrow{\scriptstyle\pi_{r'}}\\[-1mm]
 \Gamma&\xrightarrow{\ \Phi_{m|s}\ }&\Gamma'
 \end{array}
 \label{eq:state-square}
\end{equation}
Equivalently, $\pi_{r'}\circ f_{m|s}=\Phi_{m|s}\circ\pi_r$; differentiation gives the
commuting Jacobian relation
\begin{equation}
 D\pi_{r',F'}\circ Df_{m|s,F}
 =D\Phi_{m|s,\Gamma}\circ D\pi_{r,F}.
 \label{eq:commuting}
\end{equation}
Vertical variations $\delta F=F\Omega$, $\Omega^\dagger=-\Omega$, lie in
$\ker D\pi_{r,F}$.  A horizontal variation satisfies $F^\dagger\delta F=0$ and
\begin{align}
 D\pi_{r,F}[\delta F]
 &=2(\delta F F^\dagger+F\delta F^\dagger),\nonumber\\
 \|D\pi_{r,F}[\delta F]\|_F^2&=8\|\delta F\|_F^2.
 \label{eq:isometry}
\end{align}
The same metric factor appears at input and output, including for rectangular,
rank-changing branches.  Hence the horizontal quotient derivative
$\mathcal H_{F'}Df_{m|s,F}|_{\mathcal H_F}$ and the covariance derivative restricted to
the pure-state tangent have identical singular values.  A raw derivative of Householder
or QR representatives also contains vertical gauge conventions and pivot choices and is
not the invariant object in this statement.

For completeness, the one-leg factors also intertwine the initial and final pure
covariances:
\begin{equation}
 \boxed{\Gamma' L_m=L_m\Gamma.}
 \label{eq:intertwine}
\end{equation}
To verify this directly, put $\eta=\eta_m$, $d=1+\eta g$, and use the
$P_\chi\oplus Q_\chi$ decomposition
\begin{align}
 \Gamma&=\begin{pmatrix}g&w^\dagger\\w&B\end{pmatrix},\nonumber\\
 Bw&=-gw,\qquad w^\dagger w=1-g^2.
 \label{eq:pure-block}
\end{align}
Purity $\Gamma^2=\Id$ gives the two middle relations.  Equations
\eqref{eq:mobius}--\eqref{eq:L} become
\begin{align}
 L_m&=\begin{pmatrix}0&0\\-\eta w/d&\Id\end{pmatrix},\nonumber\\
 \Gamma'&=\begin{pmatrix}\eta_s&0\\0&B-\eta ww^\dagger/d\end{pmatrix}.
 \label{eq:block-L}
\end{align}
Because $(B-\eta ww^\dagger/d)w=-\eta w$, multiplication proves
Eq.~\eqref{eq:intertwine}.  Thus $L_m$ maps occupied and empty eigenspaces of
$\Gamma$ into the corresponding eigenspaces of $\Gamma'$, including exact nulls.
The intertwining identity uses $\Gamma^2=\Id$ and need not hold for a mixed state.

\section{G5: pure-state physical tangent spectrum}

At a pure state, a physical Hermitian tangent obeys $X^\dagger=X$ and
$\{\Gamma,X\}=0$.  If $U_{\Occ}$ and $U_{\Emp}$ are complete occupied and
empty frames, every such tangent is uniquely parameterized by a complex matrix $K$:
\begin{equation}
 X=2\bigl(U_{\Occ}KU_{\Emp}^\dagger
          +U_{\Emp}K^\dagger U_{\Occ}^\dagger\bigr).
 \label{eq:pure-tangent}
\end{equation}
Set $B_b=A_{T:0}U_b$ for $b\in\{\Occ,\Emp\}$.  The intertwining identity keeps
the two images in orthogonal final eigenspaces, and
\begin{equation}
 X_T=2\bigl(B_{\Occ}KB_{\Emp}^\dagger+\mathrm{h.c.}\bigr).
 \label{eq:evolved-tangent}
\end{equation}
With singular-value decompositions
$B_b=W_b\Sigma_bR_b^\dagger$, vectorization gives
\begin{equation}
 \operatorname{vec}(B_{\Occ}KB_{\Emp}^\dagger)
 =(B_{\Emp}^{*}\otimes B_{\Occ})\operatorname{vec}K.
 \label{eq:vec}
\end{equation}
Therefore the physical tangent singular spectrum is the pair product
\begin{align}
 \sigma_{ij}^{\rm phys}&=\sigma_i^{\Occ}\sigma_j^{\Emp},\nonumber\\
 \lambda_{ij}^{\rm phys}
 &=\frac{\log\sigma_i^{\Occ}+\log\sigma_j^{\Emp}}{T_{\rm obs}}.
 \label{eq:pair-rates}
\end{align}
Over the real Hermitian tangent space, the real and imaginary parts of each $K_{ij}$
give two equal real singular directions; G5 stores the complex pair once.

The numerical G5 experiment implements the pair-product estimator~\eqref{eq:pair-rates} without forming an
$N^2\times N^2$ superoperator.  Each trajectory begins from an independently seeded
Haar-random, half-filled Slater frame and is evolved natively with batch size one.  The
tangent initializer forms and diagonalizes the restricted $FF^\dagger$ once at cycle zero
to build complete occupied and empty bases; subsequent live updates are analytic and
matrix-free.  It QR-stabilizes the two blocks
separately after every cycle while multiplying the corresponding triangular factors into
scaled cumulative cores.  Cycles $1,\ldots,N_y$ align the frames; at $N_y$ only the cores
and log accumulators are reset, not the aligned bases.  At the declared integer-rounded
quarter-, half-, three-quarter-, and full-window checkpoints, direct SVDs give cumulative
core singular values.  All occupied--empty sums are formed before the
16 rates closest to zero are selected.  Core values below threshold enter as $-\infty$;
only finite pair rates are eligible.  The observer records branch/Born minima and QR
null masks; invalid branches deactivate, while core-SVD residual and physical-pair Gram
tolerances are hard acceptance gates.  Per-cycle QR diagonals are diagnostics, not
substitutes for the finite-window core SVD.  Exact null masks are legitimate, not failures.
After an exact null, however, a QR null-space completion supports the fresh-window
physical interpretation only if the preserved blocks remain complete current eigenspace
bases; a rigorous restart would rebuild them from the current frame.
Production uses $N_x=20$, $N_y=20,30,40,50,60$, $T=2N_y$, and 25 independent Born
trajectories for each wall and matched-trivial case.  No $\pm\epsilon$ trajectories,
replayable record, Choi state, or Born-score tangent is generated.

\section{G4: max-mix trajectory-operator spectrum}

G4 extracts a different object.  Start from
$\hat\rho_0=\hat{\mathbf{1}}/2^{N_{\rm orb}}$ and let
$\hat V_{\boldsymbol m}$ be the many-body evolution operator (mEO) for one realized
record.  The normalized endpoint and its weight are
\begin{equation}
 \hat\rho_{\boldsymbol m}
 =\frac{\hat V_{\boldsymbol m}\hat V_{\boldsymbol m}^\dagger}{Z_{\boldsymbol m}},
 \qquad
 Z_{\boldsymbol m}=2^{N_{\rm orb}}p_{\boldsymbol m}.
 \label{eq:maxmix-state}
\end{equation}
If $\nu_i$ are its natural occupations, Gaussian factorization gives every normalized
many-body eigenvalue and absolute mEO singular value:
\begin{align}
 \lambda_{\boldsymbol n}
 &=\prod_i\nu_i^{n_i}(1-\nu_i)^{1-n_i},\nonumber\\
 \sigma_{\boldsymbol n}(\hat V_{\boldsymbol m})
 &=\sqrt{Z_{\boldsymbol m}\lambda_{\boldsymbol n}}.
 \label{eq:maxmix-spectrum}
\end{align}
Define
\begin{align}
 e_i&=\log\frac{1-\nu_i}{\nu_i},&n_i^{(0)}&=\mathbf1_{\nu_i>1/2},\nonumber\\
 q_i&=1-2n_i^{(0)},&\delta_i&=\frac{|e_i|}{2}.
 \label{eq:operator-costs}
\end{align}
For a flipped subset $\mathcal S$, the tower consists of sums $\sum_{i\in\mathcal S}\delta_i$
sorted at fixed charge $q(\mathcal S)=\sum_{i\in\mathcal S}q_i$.  Its leading level is
\begin{align}
 \log\sigma_0=\frac12\biggl[&N_{\rm orb}\log2+\log p_{\boldsymbol m}\nonumber\\
 &+\sum_i\log\max(\nu_i,1-\nu_i)\biggr].
 \label{eq:leading-sigma}
\end{align}
Exact caps $\nu_i=0,1$ have infinite flip cost and retain an occupied/empty orientation.

Numerically, G4 evolves centered max-mix covariances in GPU batches of at most five.  The
streaming observer records convergence and impurity summaries every cycle but performs a
full natural-occupation eigendecomposition only at declared early, logarithmic, and
$N_y,3N_y/2,2N_y$ checkpoints.  It stores per-cycle conditional log-probability
increments, cumulative sums, and event counts; the cumulative $\log p_{\boldsymbol m}$ in
Eq.~\eqref{eq:maxmix-state} supplies $Z_{\boldsymbol m}$.  The occupations define the full
factorized Fock spectrum, but the implementation saves its leading scale and a low
charge-resolved subset-sum tower, not an explicit $2^{N_{\rm orb}}$ list.  Eigenpair and
natural-orbital Gram residuals are acceptance gates.  G4 uses the same sizes, duration,
and 25-trajectory wall/control counts as G5, but saves no covariance history, tangent
cocycle, Choi state, or replayable ordered-outcome record.

For a fixed max-mix record there is also an exact endpoint dictionary for the normalized
half-Jacobian $A$:
\begin{align}
 AA^\dagger&=4C_T(\Id-C_T)=\Id-\Gamma_T^2,\nonumber\\
 j_i&=2\sqrt{\nu_i(1-\nu_i)}=\operatorname{sech}\delta_i.
 \label{eq:maxmix-dictionary}
\end{align}
It explains why max-mix occupations determine tangent magnitudes, but it does not equate
G4 with G5: their initial states, Born measures, seeds, records, normalizations, and
reported spectra differ.

\section{Reference ancillas and spacetime anisotropy}

The two-dimensional topological slab supplies one-dimensional domain-wall boundaries,
whose low-energy theory is a $1+1$-dimensional CFT: one edge coordinate $y$ plus time.
Here $L_y$ is the circumference of one boundary.  The transverse size $N_x$ controls
slab width and edge--edge coupling, not the CFT circumference; the two domain walls carry
separate, oppositely chiral edge copies.  For dynamical exponent $z=1$, the expected
scaling is $t_\star\propto L_y$.  For a primary field of dimension $\Delta$, write
$z'=\alpha_{\rm st}t+\ii y$.  The cylinder correlators have spatial and temporal factors
$\sin(\pi\Delta y/L_y)$ and $\sinh(\pi\alpha_{\rm st}t/L_y)$:
\begin{align}
 g_{\rm sp}(\Delta y,0)&=\left(\frac{\pi}{L_y}\right)^{2\Delta}
 \left[\sin\left(\frac{\pi\Delta y}{L_y}\right)\right]^{-2\Delta},
 \label{eq:g-space}\\
 g_{\rm tm}(0,t)&=\left(\frac{\pi}{L_y}\right)^{2\Delta}
 \left[\sinh\left(\frac{\pi\alpha_{\rm st}t}{L_y}\right)\right]^{-2\Delta}.
 \label{eq:g-time}
\end{align}
At $\Delta y=L_y/2$, the sine is unity.  Matching the spatial and temporal factors
cancels $\Delta$ and gives~\cite{Zabalo2022}
\begin{equation}
 \boxed{\alpha_{\rm st}
 =\frac{\operatorname{arsinh}(1)L_y}{\pi t_\star}
 =\frac{\log(1+\sqrt2)L_y}{\pi t_\star}.}
 \label{eq:alpha-st}
\end{equation}
The subscript distinguishes this spacetime conversion from the circuit parameters
$\alpha_1$ and $\alpha_2$.

The numerical proxy is mutual information between two spectator references.  For each of
the two canonical modes of a wall cell, first measure its occupation $m$, append a
reference mode in occupation $1-m$, and apply the number-conserving beam splitter
\begin{equation}
 \begin{pmatrix}\hat f'\\\hat r'\end{pmatrix}
 =\frac1{\sqrt2}\begin{pmatrix}1&1\\-1&1\end{pmatrix}
 \begin{pmatrix}\hat f\\\hat r\end{pmatrix}.
 \label{eq:beam-splitter}
\end{equation}
Each pair contains one fermion and has reference entropy $\log2$ for either outcome.  One
reference subsystem therefore contains two modes and has fresh entropy $2\log2$; two
insertions append four spectator rows to the live occupied frame.  Subsequent controller
measurements have support only on the original physical rows, but physical--reference
correlations continue to evolve through the common frame.  If $\zeta_a$ are eigenvalues
of a reduced reference correlation matrix,
\begin{align}
 S(A)&=-\sum_a\bigl[\zeta_a\log\zeta_a
 +(1-\zeta_a)\log(1-\zeta_a)\bigr],\nonumber\\
 I(R_1{:}R_2)&=S(R_1)+S(R_2)-S(R_1R_2).
 \label{eq:reference-mi}
\end{align}
Only $2\times2$ and $4\times4$ matrices are diagonalized.
The late-time matching point is defined operationally, before any curve fit, by
\begin{equation}
 I_{\rm tm}(t_\star)=I_{\rm sp}(L_y/2).
 \label{eq:mi-match}
\end{equation}
It asks how many circuit cycles give the same effective separation as halfway around the
edge.  Equation~\eqref{eq:alpha-st} converts this point under the CFT cylinder hypothesis;
the matching equality itself is not evidence for CFT behavior.

The $L_y=24$ CPU audit evaluates the matching estimator directly, not by fitting
Eqs.~\eqref{eq:g-space}--\eqref{eq:g-time}.  For each $m=2,3,4,6$, 32 complete Born
trajectories are equilibrated to $mL_y$.  From each reused checkpoint, $R_1$ is inserted
at a sampled wall cell.  The spatial branch inserts $R_2$ simultaneously at separation
$L_y/2=12$; the temporal branches insert it at the same cell after
$\Delta t\in\{1,3,6,9,12,15,18\}$.  After the second insertion the same branch is followed
for $3L_y=72$ cycles.  Relative cycles 49--72, the final $L_y=24$ values, are averaged
first within each trajectory and then arithmetically over the 32 trajectories.  The
estimator requires the shortest temporal point to exceed the spatial mean and takes the
first adjacent downward crossing on the ordered sampled grid, linearly interpolated;
global monotonicity is not assumed.

For example, at equilibration $2L_y$, the saved coarse means are
$I_{\rm sp}=0.32858$, $I_{\rm tm}(1)=0.51091$, and
$I_{\rm tm}(3)=0.20135$.  Hence
\begin{align}
 t_\star&=1+\frac{0.51091-0.32858}{0.51091-0.20135}(3-1)\nonumber\\
 &=2.178.
 \label{eq:worked-crossing}
\end{align}
Applying the identical operation to all completed coarse checkpoint files gives
Table~\ref{tab:coarse}.

\begin{table}[t]
\caption{Worked coarse-audit point estimates.  These arithmetic means and interpolations
are an estimator example, not calibrated values of $t_\star$.}
\label{tab:coarse}
\centering
\footnotesize
\setlength{\tabcolsep}{3.2pt}
\begin{tabular}{@{}cccc@{}}
\toprule
equil. & $I_{\rm sp}$ & temporal bracket & interpolated $t_\star$\\
\midrule
$2L_y$ & 0.32858 & $I(1)=0.51091$, $I(3)=0.20135$ & 2.178\\
$3L_y$ & 0.12656 & $I(6)=0.19559$, $I(9)=0.07446$ & 7.710\\
$4L_y$ & 0.28096 & $I(1)=0.44209$, $I(3)=0.20006$ & 2.331\\
$6L_y$ & 0.22691 & $I(1)=0.47091$, $I(3)=0.18730$ & 2.721\\
\bottomrule
\end{tabular}
\end{table}

The $3L_y$ row illustrates sensitivity to the measured spatial baseline: its unusually
low $I_{\rm sp}$ moves the first crossing from the $1$--$3$ interval to $6$--$9$.
Table~\ref{tab:coarse} excludes refinement points $\Delta t=2,7,8$ and the paired
trajectory bootstrap and convergence gates.  It therefore carries no uncertainty
interval and is not converted into a reported $\alpha_{\rm st}$.
Evidence for the CFT interpretation would additionally require collapse of full spatial
and temporal curves under this rescaling at multiple $L_y$.

\section{Distinctions}

The three calculations answer different questions.  G4 uses a maximally mixed input to
reconstruct the absolute singular spectrum of a realized mEO, including its normalization
and charge-resolved many-body tower.  G5 uses independent pure trajectories to estimate
the Grassmannian Jacobian of the normalized covariance dynamics over a late finite-time
window.  The reference calculation instead compares a state observable at two spacetime
separations to calibrate the conversion between circuit cycles and the spatial cylinder.
The frame/M\"obius equivalence relates two coordinate descriptions of G5; neither it nor
the max-mix endpoint identity collapses these three estimators into one object.

\clearpage
\bibliographystyle{apsrev4-2}
\bibliography{references}

\end{document}
